इतरस्तु \bht{तथापि भवन्ती}त्यनेन चेत्स्वरूपप्रतिष्ठत्वमुच्यते, \bht{न तथे}ति च स्वरूपप्रतिष्ठत्वनिषेधश्चेत्, एकस्य वस्तुनः तथाभावश्चातथाभावश्च विरुध्यत इति वीक्षापन्नः पृच्छति---\bht{कथन्तर्ही}ति। इतर अह---

\str{वृत्तिसारूप्यमितरत्र॥४॥}

इति॥

कस्मात्पुनर्वृत्तिसारूप्यम? \bht{दर्शितविषयत्वात्}। उभयत्र स्वरूपप्रतिष्ठत्वाविशेषे ऽपि वृत्तिसारूप्यासारूप्यकृतो विशष इति न तथाभवातथाभावविरोधः॥

ननु च वृत्तिसारूप्ययोगे सत्यवस्थान्तरयोगात्परिणामितवादिदोषः प्राप्नोत्येव॥

न। दर्शितविषयत्वेन परिहृतत्वात्। चित्तवृत्यध्यारोपितं हि तन्न स्वतः, स्फटिकाद्युपधानोपरागवत्॥

\bht{व्युत्थान} इति निरोधादन्यत्र। \bht{याश्चित्त}स्य क्लिष्टादयो \bht{वृत्तयः}, \bht{ता}भिर\bht{विशिष्टा} सदृशी \bht{वृत्ति} स्वरूपमस्य। % check if there was a visarga after v.rtti
चित्तवृत्तिव्यतिरेकेण स्वरूपव्यतिरेकेण च पुरुषस्य वृत्यभावात्। सदाज्ञातविषयत्वे परिणामानुपपत्तेश्च॥

\bht{तथा च} पूर्व्वाचार्य्य\bht{सूत्रम्}--\bht{एकमेव दर्शनम्} बुद्धिपुरुषयोः। किन्तदित्याह---\bht{ख्यातिरेव दर्शनम्}। बुद्धिवृत्तिरेव दर्शनम्। ख्यायते पुरुषेणेति च ख्यातिः। ख्यायते ऽनया च बुद्धिपुरुषयोः स्वरूपमिति ख्यातिः। सैवार्त्थाकारतया गृह्यत इति करणम्। स्वेन च ख्यातिरूपेण गृह्यत इति कर्म्म च भवति। तथा दृश्यत इति दर्शनम्। दृश्यते ऽनेनेति च ज्ञायते ऽनेनेति च सर्व्वमेवमादि द्रष्टव्यम्॥

यत्तु ख्यातिकर्तृत्वे चित्तस्य तदेव ज्ञायते जानाति चेति ज्ञानान्तरकल्पनावैयर्त्थ्यप्रसंग इति तस्य `न तत्स्वाभासन्दृश्यत्वादि'त्यत्रैव परिहारं वक्ष्यामः॥ % Is this sentence coherent? Yes, it is. It refers to YS 4.19.

ननु च यदि चित्तम्परिविजृम्बहितमेवेदं सर्व्वम्पुरुषश्चेदुदास्ते कथमसौ भोक्तेति—न, कारकवैचित्र्यात्। कानिचिद्व्याप्रियमाणानि कारकाणि कानिचिद्व्यापारादृते सन्निधिमात्रेण कर्तृत्वं प्रतिपद्यन्ते। यथा पचतीत्यधिश्रयणादिषु व्याप्रियमाणः पुरुषः कर्तेत्युच्यते। स्थाली सम्भवधारणाभ्यामव्याप्रियमाणैव पचतित्युच्यते। तथा स्थालीयमध्यगतमाकाशमवकाशन्ददाति। न तस्य व्यापारं प्रतीमः। यथा च राजादर्शनमात्रेणैव कर्त्ता। सर्व्वाणि च कारकाणि स्वेषु व्यापारेषु कर्त्तॄणि।





[14,10]	itaras tu tathaa^ api bhavanty anena cet svaruupa-prati.s.thatvam ucyate 
na tathaa^ iti ca svaruupa-prati.s.thatvani.sedha"s ced ekasya vastunas tathaabhaava"s 
ca^ atathaa-bhaava"s ca virudhyata^ iti viik.saapanna.h p.rcchati katha.m tarhi 
iti./ itara^ aaha v.rttisaaruupyam itaratra^ iti./ kasmaat punar v.rttisaaruupyam ? 
dar"sitavi.sayatvaat // ubhayatra svaruupa-prati.s.thatvaavi"se.se^ api 
v.rtti-saaruupyaasaaruupyak.rto vi"se.sa^ iti, na tathaabhaavaatathaabhaava-virodha.h./

[14,15]	nanu ca v.rtti-saaruupya-yoge sati^ avasthaantara-yogaat 
pari.naamitva-aadi-do.sa.h praapnoty eva. na. dar"sitavi.sayatvena parih.rtatvaat. 
citta-v.rtty-adhyaaropita.m hi tat, na svata.h 
spha.tika-aady-upadhaana-uparogavat. vyutthaana^ iti./ nirodhaad anyatra ya"s 
cittasya kli.s.ta-aadayo v.rttaya.h taabhir avi"si.s.taa sad.r"sii v.rtti.h svaruupam 
asya, citta-v.rtti-vyatireke.na svaruupa-vtyatireke.na ca puru.sasya 
v.rtty-abhaavaat sadaa j~naata-vi.sayatve pari.naama-anupapatte"s ca./

[14,21]	tathaa ca puurva-aacaarya-suutram ekam eva dar"sana.m buddhi-puru.sayo.h. 
ki.m tad ity aaha khyaatir eva dar"sanam. buddhi-v.rttir eva dar"sanam. khyaayate 
puru.se.na^ iti ca khyaati.h khyaate 'nayaa ca buddhi-puru.sayo.h svaruupam^ iti 
khyaati.h. saa^ eva-arthaakaaratayaa [arthaakaaro 'nayaa] g.rhyata^ iti kara.na.m svena 
ca khyaati-ruupe.na g.rhyata^ iti karma ca bhavati. tathaa d.r"syata^ iti dar"sana.m, 
d.r"syate 'neneti ca [j~naayata^ iti j~naana.m ] j~naayate 'neneti ca sarvam 
evam-aadi dr.s.tavyam./

[14,27]	yat tu khyaati-kart.rtve cittasya tad eva j~naayate jaanaati ca^ iti 
(j~naanaad iti ca^ iti) j~naanaantarakalpnaa-vaiyarthya-prasa"nga^ iti tasya 'na 
tat-svaabhaasa.m d.r"syatvaat' ity atraiva parihaara.m vak.syaama.h./

[15,4]	nanu ca yadi citta-parvij.rmbhitam eveda.m sarva.m paru.sa"s ced udaaste 
katham asau bhokteti ? na kaarakavaicitryaat. kaani-cid vyaapriyamaa.naani 
kaarakaa.ni, kaanicid vyaapaaraad .rte sannidhi-maatre.na kart.rtva.m pratipadyante./

[15,7]	yathaa pacati^ ity adhi"sraya.na-aadi.su vyaapriyamaa.na.h puru.sa.h kartety 
ucyate. sthaalii sambhava-dhaara.naabhyaam avyaapriyam.naiva pacati^ ity ucyate. 
tathaa sthaalii-madhya-gatam aakaacam avakaaca.m dadaati, na tasya vyaapaara.m 
pratiima.h./

[15,10]	yathaa ca raajaa dar"sana-maatre.na^ eva kartaa. sarvaa.ni ca kaarakaa.ni 
sve.su vyaapaare.su kartY.ni. na ca^ aaditya.h prakaa"sayan kara.naantaram apek.sate 
vyaapriyate vaa / na hi prakaa"sanam avidyamaanam utpaadayati gamana-aadivat. 
prakaa"saatmataa^ eva hi tasya sad-bhaava.h tathaapi d.r"si-maatra.na gha.ta-aadiinaa.m 
prakaaca-ruupe.na^ abhivyakte.h prakaacayati^ ity ucyate, evam iha^ api 
d.r"si-maatre.na puru.se.na d.r"syamaanaa.m citta-v.rttiinaa.m cid-aatmanaa vyaapyamaanatvaat 
dra.s.taa puru.sa tiimam artha.m dar"sayati cittam ayaskaantama.ni-kalpam iti./

[15,16]	yasya^ apy aatma-mans-sa.myogaat j~naanam utpadyate tasya^ api j~naanena 
j~neya.m vyaapnuvan naatmaa j~naatety ucyate, na vyaapriyamaana.h. j~naanasya^ api 
[j~naanam api] gu.natve sati ni.skriyatvaad aatmana^ eva^ artha.m dar"sayati. na^ 
api kriyaan abhyupagamaat jaanaati^ ity-aadi na^ upapadyate. j~naanena hi j~neya.m 
vyaapnivan jaanaati^ ity ucyate./

[15,20]	ki.m ca^ anyat aatma-manas-sa.myogaad apuurva-j~naanotpattau saty 
aatma-manas-taj-j~naana.m na manasa.h. aatmaa^ eva j~nnaati na mana^ iti ko hetur 
bhavet ? atha^ api syaat aatmana.h samavaayi-kaara.natvaad iti. na^ etad eva.m 
asiddhatvaat. yathaa^ eva j~naanam aatmana^ iti saadhya.m evam aaaatmana^ eva 
samavaayi-kaara.nattva.m na manasa^ iti saadhyam anaatma-sva-ruupatve saty 
ubahaya-sa.myogajatva-abhyupagamaat./

[15,25]	atha^ api syaat j~naanaahita-sa.mskaarasyaatma[na.h]sm.rty-ekakart.rtvena 
pratisandhaana-dar"sanaad iti na tasya^ api saadhyatvaat. tatra^ api 
j~naana-sa.mskaara-sm.rti-pratisandhaanaadiiny\label{JSSP04} aatmana^ eva na manasa^ iti 
saadhaniiyaani. phalena samaanaa"srayatvaad icchaa-aadiinaam iti cet na tatra^ api 
tulyatvaat. sa.myogajatve saty aatmana.h phalecchaadiini na manasa^ iti du.hsaadham 
eva./

[16,4]	ki.m ca^ asya j~naanasya^ api j~neyatva-abhyupagamaad 
anavasthaa-prasa"nga.h. na hi j~naanam anyena kara.na-bhuutena j~naanena vij~naatu.m 
"sakya.m tad apy anyena tad apy anyenety ani.s.tha-prasa"ngaat./

[16,6]	athaanavasthaa-prasa"nga-parihaaraaya duuram api gatvaa ki.mcij j~naanam 
aadya.m vaa j~naanaantare.na na g.rhyate svayam aatmanaa^ eva g.rhyata^ iti tathaa 
j~naanavad arthasya^ api svaya.m graha.na-prasa"nga.h./

[16,8]	atha diipavat svayam avabhaasate ca^ avabhaasayati ca 
avabhaasaatmakatvaat arthasyaanavabhaasaatmakatvaad ado.sa^ iti cet tathaapi 
prakaa"sa-svaruupe.sv arthe.su vij~naanenaagraha.na-prasa"nga.h./

[16,11]	atha jjaanasya^ aatma-samavetatvaavabhaasana-ruupaabhyaa.m svaya.m graha.na.m 
bahi.hprakaa"saatmakaanaa.m tu j~naana-vipariita-dharnatvaad ado.sa^ iti cet tac ca na 
yadi ca^ aatma-samavetatva.m prakaa"saatmakatva.m ca j~naanasya svaya.mgraha.ne 
hetu.h syaat sarve sarvajjaa.h praapnuvanti./

[16,14]	atha graahyaakaaratvam api kaara.na.m cet tac ca na^ anyasya 
j~naanasyaanyaakaaratayaa kara.natva.m yathaa tathaa j~naanasya^ api 
graahyasyaanyaakaare.naanyad eva j~naana.m kara.na.m syaat tathaa ca^ anavasthaa 
du.spariharaa^ eva./

[16,17]	ki.mcaanyat na hi ka"scit svaya.mgraahya-graahakatva-prasaadhanaaya 
d.r.s.taanto vidyate pradiipasya^ api cak.sur-aadi-graahyatvaat. na ca^ apy 
a.m"sa-dvayam asti j~naanasya niravayavatvaat. saty api ca^ a.m"sa-dvaye tayo.h 
prakaa"saatmakatvaan naanyonya-graahya-graahaka-bhaava.h pradiipayor iva dvayo.h./

[16,21]	ki.m ca^ anyat j~naanasya graahya-graahakatva-abhyupagame saty 
aatma-sad-bhaavo api dussampaada^ eva. tasmaad 
arthaakaara-j~naana-sm.rti-pratisandhaana-prayotnecchaa-aadiinaam anaatma-dharmatva.m 
graahyatvaat ruupa-aadiinaa.m iva. paraarthatvaac ca sa"nghaat-aa"srayatva.m 
ruupa-aadi-nidar"sanena^ eva. tathaa^ aa"srayavattvaad anityatvaac ca 
prayatra-puurvakatvaac ca^ ity evam-aadibhyo^ anaatma-dharmatva-siddhi.h./

[16,25]	ayaskaanata-ma.ni-kalpam./ ayakaantama.nir eva. yathaa^ ayakaantama.nir 
ayasa.h svaatma-vyaapti.m pratyupakaroti sannidhimaatra.na tathaa cittam api 
cit-svaruupasyaatmana.h svaatma-vyaapti.m pratyupakaroti d.r"syatvena./

[17,9]	nartakiiva dar"sita-vi.sayatvaad anta.hkara.nenety uktam tatra 
dar"sayit.rtva.m cittasya dra.s.t.rtva.m ca puru.sasya ki.mk.rtam ? 
svasvaamitvak.rtam. katha.m puna.h svasvaamitvam ? vastu-svabhaavatvaat. katham 
it ced ukta.m cittam ayaskaantama.nikalpam iti./ atha citta-v.rttiinaa.m bodhe 
ko hetu.h puru.sasya^ iti ? ucyate yasmaad eva.m sannidhi-maatra-upakaari citta.m 
d.r"syatvena tasmaac citta-v.rttiinaa.m bodhe puru.sasya^ aya.m hetu.h bhavati yo 
'sau^ anaadi-sambanda.h./ 4./
